\begin{thebibliography}{24}
\bibitem{aharonovich} Aharonovich I., Englund D., Toth M. Solid-state single-photon emitters. \textit{Nature Photonics} \textbf{10}, 631--641 (2016). \href{https://doi.org/10.1038/nphoton.2016.186}{doi:10.1038/nphoton.2016.186}

\bibitem{senellart} Senellart P., Solomon G., White A. High-performance semiconductor quantum-dot single-photon sources. \textit{Nature Nanotechnology} \textbf{12}, 1026--1039 (2017). \href{https://doi.org/10.1038/nnano.2017.218}{doi:10.1038/nnano.2017.218}

\bibitem{purcell} Purcell E. M. Spontaneous emission probabilities at radio frequencies. \textit{Physical Review} \textbf{69}, 681 (1946). \href{https://doi.org/10.1103/PhysRev.69.674.2}{doi:10.1103/PhysRev.69.674.2}

\bibitem{novotny} Novotny L., Hecht B. \textit{Principles of Nano-Optics}, 2nd ed. Cambridge University Press, Cambridge (2012). \href{https://doi.org/10.1017/CBO9780511794193}{doi:10.1017/CBO9780511794193}

\bibitem{koenderink} Koenderink A. F. Single-photon nanoantennas. \textit{ACS Photonics} \textbf{4}, 710--722 (2017). \href{https://doi.org/10.1021/acsphotonics.7b00061}{doi:10.1021/acsphotonics.7b00061}

\bibitem{akselrod} Akselrod G. M., Argyropoulos C., Hoang T. B., Ciracì C., Fang C., Huang J., Smith D. R., Mikkelsen M. H. Probing the mechanisms of large Purcell enhancement in plasmonic nanoantennas. \textit{Nature Photonics} \textbf{8}, 835--840 (2014). \href{https://doi.org/10.1038/nphoton.2014.228}{doi:10.1038/nphoton.2014.228}

\bibitem{tame2013} Tame M. S., McEnery K. R., Özdemir Ş. K., Lee J., Maier S. A., Kim M. S. Quantum plasmonics. \textit{Nature Physics} \textbf{9}, 329--340 (2013). \href{https://doi.org/10.1038/nphys2615}{doi:10.1038/nphys2615}

\bibitem{grange2015} Grange T., Hornecker G., Hunger D., Poizat J.-P., Gérard J.-M., Senellart P., Auffèves A. Cavity-funneled generation of indistinguishable single photons from strongly dissipative quantum emitters. \textit{Physical Review Letters} \textbf{114}, 193601 (2015). \href{https://doi.org/10.1103/PhysRevLett.114.193601}{doi:10.1103/PhysRevLett.114.193601}

\bibitem{anger} Anger P., Bharadwaj P., Novotny L. Enhancement and quenching of single-molecule fluorescence. \textit{Physical Review Letters} \textbf{96}, 113002 (2006). \href{https://doi.org/10.1103/PhysRevLett.96.113002}{doi:10.1103/PhysRevLett.96.113002}

\bibitem{kuhn} Kühn S., Håkanson U., Rogobete L., Sandoghdar V. Enhancement of single-molecule fluorescence using a gold nanoparticle as an optical nanoantenna. \textit{Physical Review Letters} \textbf{97}, 017402 (2006). \href{https://doi.org/10.1103/PhysRevLett.97.017402}{doi:10.1103/PhysRevLett.97.017402}

\bibitem{torma2015} Törmä P., Barnes W. L. Strong coupling between surface plasmon polaritons and emitters: a review. \textit{Reports on Progress in Physics} \textbf{78}, 013901 (2015). \href{https://doi.org/10.1088/0034-4885/78/1/013901}{doi:10.1088/0034-4885/78/1/013901}

\bibitem{chikkaraddy2016} Chikkaraddy R., de Nijs B., Benz F., Barrow S. J., Scherman O. A., Rosta E., Demetriadou A., Fox P., Hess O., Baumberg J. J. Single-molecule strong coupling at room temperature in plasmonic nanocavities. \textit{Nature} \textbf{535}, 127--130 (2016). \href{https://doi.org/10.1038/nature17974}{doi:10.1038/nature17974}

\bibitem{fischer2016} Fischer K. A., Müller K., Lagoudakis K. G., Vučković J. Dynamical modeling of pulsed two-photon interference. \textit{New Journal of Physics} \textbf{18}, 113053 (2016). \href{https://doi.org/10.1088/1367-2630/18/11/113053}{doi:10.1088/1367-2630/18/11/113053}

\bibitem{ge2014} Ge R.-C., Hughes S. Design of an efficient single photon source from a metallic nanorod dimer: a quasi-normal mode finite-difference time-domain approach. \textit{Optics Letters} \textbf{39}, 4235--4238 (2014). \href{https://doi.org/10.1364/OL.39.004235}{doi:10.1364/OL.39.004235}

\bibitem{franke2019} Franke S., Hughes S., Kamandar Dezfouli M., Kristensen P. T., Busch K., Knorr A., Richter M. Quantization of quasinormal modes for open cavities and plasmonic cavity quantum electrodynamics. \textit{Physical Review Letters} \textbf{122}, 213901 (2019). \href{https://doi.org/10.1103/PhysRevLett.122.213901}{doi:10.1103/PhysRevLett.122.213901}

\bibitem{hughes2019} Hughes S., Franke S., Gustin C., Kamandar Dezfouli M., Knorr A., Richter M. Theory and limits of on-demand single-photon sources using plasmonic resonators: a quantized quasinormal mode approach. \textit{ACS Photonics} \textbf{6}, 2168--2180 (2019). \href{https://doi.org/10.1021/acsphotonics.9b00849}{doi:10.1021/acsphotonics.9b00849}

\bibitem{zhang2025} Zhang S., Traverso A., Dolgopolova E., Singh A., Kishida H., Livshits M., Sheehan C., Bowes E., Li C., Hollingsworth J., Mikkelsen M. Solution-processed ultrafast, room-temperature single-photon source at 1550 nm. \textit{ACS Nano} \textbf{19}, 19035--19045 (2025). \href{https://doi.org/10.1021/acsnano.4c18261}{doi:10.1021/acsnano.4c18261}

\bibitem{khatua} Khatua S., Paulo P. M. R., Yuan H., Gupta A., Zijlstra P., Orrit M. Resonant plasmonic enhancement of single-molecule fluorescence by individual gold nanorods. \textit{ACS Nano} \textbf{8}, 4440--4449 (2014). \href{https://doi.org/10.1021/nn406434y}{doi:10.1021/nn406434y}

\bibitem{mohammadi} Mohammadi A., Sandoghdar V., Agio M. Gold nanorods and nanospheroids for enhancing spontaneous emission. \textit{New Journal of Physics} \textbf{10}, 105015 (2008). \href{https://doi.org/10.1088/1367-2630/10/10/105015}{doi:10.1088/1367-2630/10/10/105015}

\bibitem{paper1} Sadidi A. H., Molaei Tabari S. M. P., Bagheri Harouni M. Resonant versus geometric anisotropy of the Purcell enhancement near a gold nanorod: plasmon-mode decomposition, an equal-volume control and design rules for single-photon sources. Companion manuscript, in preparation (2026).

\bibitem{johnson} Johnson P. B., Christy R. W. Optical constants of the noble metals. \textit{Physical Review B} \textbf{6}, 4370--4379 (1972). \href{https://doi.org/10.1103/PhysRevB.6.4370}{doi:10.1103/PhysRevB.6.4370}

\bibitem{hohenester2012} Hohenester U., Trügler A. MNPBEM -- A Matlab toolbox for the simulation of plasmonic nanoparticles. \textit{Computer Physics Communications} \textbf{183}, 370--381 (2012). \href{https://doi.org/10.1016/j.cpc.2011.09.009}{doi:10.1016/j.cpc.2011.09.009}

\bibitem{qutip5} Lambert N., Giguère E., Menczel P., Li B., Hopf P., Suárez G., Gali M., Lishman J., Gadhvi R., Agarwal R., Galicia A., Shammah N., Nation P., Johansson J. R., Ahmed S., Cross S., Pitchford A., Nori F. QuTiP 5: the quantum toolbox in Python. \textit{Physics Reports} \textbf{1153}, 1--62 (2026). \href{https://doi.org/10.1016/j.physrep.2025.10.001}{doi:10.1016/j.physrep.2025.10.001}

\bibitem{yorulmaz2012} Yorulmaz M., Khatua S., Zijlstra P., Gaiduk A., Orrit M. Luminescence quantum yield of single gold nanorods. \textit{Nano Letters} \textbf{12}, 4385--4391 (2012). \href{https://doi.org/10.1021/nl302196a}{doi:10.1021/nl302196a}
\end{thebibliography}
